%% kapa: exam=1 round=36 year=2025 period=1 form=A subject=economics subject_ko=경제학원론 seq=32 number=72
%% source: Q-Net 감정평가사 기출문제 (cst00309) artlSeq=5231525 file=raw/1008/5231525/2025년 제36회 감정평가사 1차시험 1교시 문제지 원본.pdf
%% conversion: kordoc 4.24.51 -> tools/gen1.py (auto); review=proofread-sonnet; flags=none
\begin{question}{72}
필립스 곡선에 관한 설명으로 옳지 않은 것은?
\begin{choices}
  \item 우하향하는 필립스 곡선에서 인플레이션 감축을 위해서는 실업률 증가를 감수해야 한다.
  \item 단기 총공급곡선이 가파를수록 단기 필립스곡선은 가파른 모양을 가진다.
  \item 물가상승률과 경제성장률의 관계에 오쿤의 법칙을 적용하면 필립스 곡선을 도출할 수 있다.
  \item 합리적 기대에 따라 필립스 곡선이 즉시 이동하면, 실업율의 증가 없이 인플레이션을 감축할 수 있다.
  \item 기대 인플레이션이 상승하면 필립스 곡선은 좌측으로 이동한다.
\end{choices}
\end{question}
\answer{5}
